%-------------------------------------------------------------------------------------------------
% 3.a) Listar até 5 pesquisas/produtos mais relevantes (incluindo pesquisas/produtos/contribuições 
% ao desenvolvimento tecnológico e à inovação dos últimos 5 anos), podendo ser artigos científicos, 
% livros, capítulos de livros, palestras a convite, patentes (solicitadas, concedidas e especialmente 
% as licenciadas), softwares registrados, cultivares, empresas de base tecnológica criadas, relatórios 
% técnicos de projetos que liderou ou em que teve papel destacado, novos produtos, novos processos ou 
% outros tipos de resultados de pesquisa mais relevantes que se relacionam a seu projeto desta submissão. 
% Todos os autores, até um máximo de 15, deverão constar da citação. Usar o padrão bibliográfico 
% convencional em sua área. Das listas com mais de 5 itens, a FAPESP somente utilizará no processo de 
% análise os 5 primeiros.

% 3.b) Incluir, em cada resultado, a justificativa da escolha: impacto (social, científico, econômico, 
% tecnológico, em inovação etc.), relevância para a análise de sua produção – científica, tecnológica 
% e em inovação – e importância para seus projetos de pesquisa em andamento/propostos, em no máximo 5 
% linhas.
%-------------------------------------------------------------------------------------------------
\begin{cvpubs}
    \cvpub{1. Wayne B, Prince D, \textbf{Tal F}, Kent C, Allen B. AAAA. Título do artigo com o nome da espécie em itálico, por exemplo \textit{Rana temporaria} Linnaeus, 1758. \textit{Journal of Gotham Zoology} 86, 126927. \textbf{Citações: X.}}

\vspace{1em}
\noindent \lipsum[1][1-4]
\vspace{1em}

    \cvpub{2. Jordan H, Curry A, \textbf{Tal F}, Queen O, Kyle S, Wayne B. AAAA. Título do segundo artigo sobre \textit{Lacerta agilis} Linnaeus, 1758. \textit{Metropolis Journal of Evolution} 144, e006. \textbf{Citações: X.}}

\vspace{1em}
\noindent \lipsum[2][1-4]
\vspace{1em}

    \cvpub{3. \textbf{Tal F}, Prince D, Kent C, Allen B. AAAA. Título do terceiro artigo, uma revisão sobre a diversidade de espécies de uma região. \textit{European Journal of Kryptonian Taxonomy}, 836, 96--130. \textbf{Citações: X.}}

\vspace{1em}
\noindent \lipsum[3][1-4]
\vspace{1em}

    \cvpub{4. Curry A, Kyle S, Jordan H, Queen O, \textbf{Tal F}, Wayne B. AAAA. Título do quarto artigo, por exemplo a descrição de um genoma ou de uma nova espécie. \textit{Genes, Genomes and Gotham}, 12, jkac270. \textbf{Citações: X.} (capa)}

\vspace{1em}
\noindent \lipsum[4][1-4]
\vspace{1em}

    \cvpub{5. \textbf{Tal F}, Wayne B, Kent C, Prince D. AAAA. Título do quinto artigo, de preferência o mais relacionado ao projeto submetido. \textit{Molecular Phylogenetics of Central City} 128, 123--146. \textbf{Citações: X.}}

\vspace{1em}
\noindent \lipsum[5][1-4]
\vspace{1em}
    
\end{cvpubs}
